\ncertchapter{2}{Radioactive Decay}\chlabel{ch:decay}

\chref{ch:nucleus} treated nuclei as fixed objects with definite masses and binding energies. Most
nuclides are not fixed at all. Of the more than $3000$ nuclides that have been made and studied, only
about $251$ are stable; every other one, sooner or later, changes into something else by itself, with no
push from outside. This chapter asks what such an unstable nucleus does: which particles it emits, how
much energy they carry away, which nuclei are unstable in the first place, and where a chain of decays
finally comes to rest.

Almost every answer comes from tools already in hand. Energy conservation (\chref{ch:energy}), extended
by $E=mc^2$ and the atomic-mass tables of \chref{ch:nucleus}, says whether a decay can happen and how much
energy it releases. Momentum conservation (\chref{ch:collisions}) says how that energy is shared among
the products, and it is momentum, more than anything else, that forced physicists to accept a new and
almost undetectable particle, the neutrino. What is genuinely new is the cast: positrons, neutrinos and
their antiparticles, which must be met before any decay equation can be written.

Two limits should be stated at the outset. This chapter is about \emph{what} comes out of a decay and
with what energy, not about \emph{when} a given nucleus decays; that question, whose answer is a law of
chance, is the subject of \chref{ch:decaylaw}. And some of the reasons behind the facts (why an $\alpha$
particle can leak out of a nucleus at all, why $\beta$ decays are slow) belong to quantum mechanics and
the theory of the weak force. They are explained here in words, and flagged where they are stated
without proof.

\section{An Accidental Discovery}\label{sec:decay-discovery}

In 1896 Henri Becquerel stored uranium salts in a drawer on top of a photographic plate wrapped in black
paper. When the plate was developed it was fogged in the outline of the crystals, although no light had
reached it. Marie and Pierre Curie then showed that the effect depended only on the \emph{amount} of
uranium present, not on the chemical compound it was in, nor on temperature or pressure. They named it
\textbf{radioactivity}, and by chemical separation they isolated two far more active elements, polonium
and radium.

The insensitivity to chemistry was the decisive clue. Chemical bonds involve energies of a few
electron-volts and affect only the outer electrons of an atom; heating a sample by a few hundred kelvin
changes the energy of each atom by only a fraction of an electron-volt. A process that ignores all of
this must originate somewhere that chemistry and heat cannot reach: the nucleus, where, as
\chref{ch:nucleus} showed, the energies are millions of electron-volts per particle. Before the emitted
radiation can be described, the particles that take part in it need names.

\section{The Cast of Particles}\label{sec:decay-particles}

Radioactivity involves a handful of particles, some of them unfamiliar. Each is defined here before it
is used anywhere else, with the mass, charge and one further property, lepton number, that the rest of
the chapter needs. Masses are given both in atomic mass units ($\SI{1}{\amu}$, one-twelfth of the mass of
a neutral carbon-12 atom) and as rest energies in MeV, linked by the conversion of \chref{ch:nucleus},
$\SI{1}{\amu}\times c^2=\SI{931.494}{\MeV}$; writing a mass as ``MeV$/c^2$'' lets $c^2$ cancel silently in
calculations.

\begin{defbox}{Proton, neutron and nucleon}
These were met in \chref{ch:nucleus}; their properties are collected here. The \textbf{proton} $p$
carries charge $+e$, where $e=\SI{1.602e-19}{C}$ is the elementary charge, and has mass
$m_p=\SI{1.007276}{\amu}=\SI{938.3}{\MeV}/c^2$. The \textbf{neutron} $n$ is electrically neutral and
slightly heavier, $m_n=\SI{1.008665}{\amu}=\SI{939.6}{\MeV}/c^2$. Protons and neutrons are the two kinds of
\textbf{nucleon}, the particles that make up a nucleus. In decay equations they are written $\nuc11{p}$
(or $\nuc11{H}$, the nucleus of hydrogen) and $\nuc10{n}$: the upper number counts nucleons and the lower
number counts charge in units of $e$.
\end{defbox}

\begin{defbox}{Electron}
The \textbf{electron} $e^-$ carries charge $-e$ and has mass $m_e=\SI{0.000549}{\amu}=\SI{0.511}{\MeV}/c^2$,
about $1/1836$ of a proton's. Electrons fill the shells of an atom around its nucleus; the innermost
shell, closest to the nucleus, is called the \textbf{K shell}. An electron emitted by a nucleus is called a
\textbf{beta-minus particle}, written $\beta^-$ or $\nuc{0}{-1}{e}$: nucleon number $0$, charge $-1$.
\end{defbox}

\begin{defbox}{Photon}
A \textbf{photon} $\gamma$ is a single quantum (smallest indivisible packet) of electromagnetic radiation.
It has no charge and no mass, always travels at the speed of light $c$, and carries energy $E=hf$, where
$f$ is the frequency of the radiation and $h=\SI{6.626e-34}{\joule\second}$ is Planck's constant, together with
momentum of magnitude $E/c$ (both stated without proof; they belong to the quantum theory of light).
Visible light, X-rays and gamma rays are all photons and differ only in energy. Photons from the nucleus
(keV to MeV) are called \textbf{gamma rays}; those from rearrangements of an atom's inner electrons (up to
about $\SI{100}{\keV}$) are \textbf{X-rays}. Photons are created and absorbed freely: there is no
conservation of photon number.
\end{defbox}

\begin{defbox}{Antiparticle}
Every particle has an \textbf{antiparticle} with exactly the same mass but the \emph{opposite} electric
charge, and the opposite value of every other conserved charge-like quantity (such as lepton number,
defined below). Antiparticles are written with a bar ($\bar\nu$) or with the opposite charge sign ($e^+$
for the antielectron). Together they make up \textbf{antimatter}. They are not rare or exotic in nuclear
physics: radioactive decays create them routinely.
\end{defbox}

\begin{defbox}{Positron}
The \textbf{positron} $e^+$ is the antiparticle of the electron: mass $\SI{0.511}{\MeV}/c^2$, charge $+e$.
Predicted by Dirac (1928) and discovered by Anderson in cosmic rays (1932), it is emitted in $\beta^+$
decay, where it is called a \textbf{beta-plus particle}, written $\beta^+$ or $\nuc{0}{+1}{e}$. A positron
does not survive long in ordinary matter. Within about a nanosecond it meets an electron and the pair
\textbf{annihilates}: both particles vanish and their total rest energy, $2m_ec^2=\SI{1.022}{\MeV}$,
reappears as photons (Figure~\ref{fig:decay-antimatter}b). For a pair essentially at rest the total momentum
is zero, so a single photon is impossible, since it would carry momentum $E/c\neq0$. Instead two photons
of $\SI{511}{\keV}$ each fly off in exactly opposite directions. The reverse process, \textbf{pair
production}, turns a photon of more than $2m_ec^2=\SI{1.022}{\MeV}$ into an electron and a positron; it
needs a nearby nucleus to take up momentum (Figure~\ref{fig:decay-antimatter}c).
\end{defbox}

\begin{defbox}{Neutrino and antineutrino}
The \textbf{electron neutrino} $\nu_e$ (Greek ``nu'') is an electrically neutral particle with an
extremely small mass (less than $\SI{1}{\eV}/c^2$, so treated here as massless) that travels at very
nearly the speed of light. It feels neither the strong nuclear force nor the electric force, only the
\textbf{weak nuclear force}, so it passes through matter almost unhindered. Its antiparticle is the
\textbf{electron antineutrino} $\bar\nu_e$ (``nu-bar''). Both are neutral and practically massless, so
they are told apart by their lepton number $L_e$ (defined below), $+1$ for $\nu_e$ and $-1$ for $\bar\nu_e$. In practice an
\emph{antineutrino} always accompanies an emitted \emph{electron} ($\beta^-$ decay), and a \emph{neutrino}
accompanies an emitted \emph{positron} ($\beta^+$ decay) or a captured electron (electron capture).
Neutrinos carry away energy and momentum but are almost never absorbed, so that energy is lost to any
detector or absorber.
\end{defbox}

\begin{defbox}{Lepton and lepton number}
A \textbf{lepton} is a particle that does not feel the strong nuclear force; in this chapter the leptons
are the electron and the electron neutrino and their antiparticles. The electron and the neutrino are
assigned \textbf{electron lepton number} $L_e=+1$; their antiparticles, the positron and the antineutrino, have
$L_e=-1$; all other particles (protons, neutrons, photons, nuclei) have $L_e=0$. (The subscript distinguishes it from the
angular momentum $L$ of \chref{ch:rotation}; heavier relatives of the electron, met only in particle
physics, carry lepton numbers of their own.) In every process ever observed the \emph{total} electron
lepton number is conserved. That is why leptons are created in balanced pairs:
$e^-$ with $\bar\nu_e$ ($+1-1=0$), or $e^+$ with $\nu_e$ ($-1+1=0$).
\end{defbox}

Figure~\ref{fig:decay-antimatter} pictures the particle--antiparticle pairs and the two processes that
convert rest energy into photons and back.

\begin{figure}[htbp]
\centering
\tikzfig{decay_antimatter}
\figcaption{fig:decay-antimatter}{(a) A particle and its antiparticle have the same mass and opposite
charge-like quantities; the neutrino and antineutrino are both neutral and differ only in lepton number
$L_e$. (b) An electron and a positron at rest annihilate into two $\SI{511}{\keV}$ photons moving in
opposite directions, which is the only way to conserve both energy and zero total momentum. (c) The
reverse: a photon above $\SI{1.022}{\MeV}$, passing close to a nucleus, can create an electron--positron
pair; the nucleus absorbs the photon's momentum.}
\end{figure}

\begin{tablebox}{The particles of radioactivity}
\centering
\renewcommand{\arraystretch}{1.15}
\begin{tabular}{@{}lcccc>{\raggedright\arraybackslash}p{5.0cm}@{}}
\toprule
\textbf{Particle} & \textbf{Symbol} & \textbf{Charge} & \textbf{Mass (MeV$/c^2$)} & $L_e$ & \textbf{Where it appears}\\
\midrule
proton & $p$, $\nuc11{p}$ & $+e$ & $938.3$ & $0$ & nuclei; target and product in reactions\\
neutron & $n$, $\nuc10{n}$ & $0$ & $939.6$ & $0$ & nuclei; free mean life about $\SI{880}{s}$\\
electron & $e^-$, $\beta^-$ & $-e$ & $0.511$ & $+1$ & atomic shells; emitted in $\beta^-$ decay; captured in EC\\
positron & $e^+$, $\beta^+$ & $+e$ & $0.511$ & $-1$ & emitted in $\beta^+$ decay; annihilates\\
neutrino & $\nu_e$ & $0$ & $\approx0$ & $+1$ & with $e^+$ in $\beta^+$ decay; in EC; from the Sun\\
antineutrino & $\bar\nu_e$ & $0$ & $\approx0$ & $-1$ & with $e^-$ in $\beta^-$ decay; from reactors\\
photon & $\gamma$ & $0$ & $0$ & $0$ & gamma decay, annihilation, X-rays\\
alpha particle & $\alpha$, $\nuc42{He}$ & $+2e$ & $3727.4$ & $0$ & a helium-4 \emph{nucleus}; emitted in $\alpha$ decay\\
\bottomrule
\end{tabular}\label{tab:decay-particles}
\end{tablebox}

\begin{remarkbox}{What is really changing in beta decay?}
Protons and neutrons are not fundamental. Each is built from three \emph{quarks}: the proton from two
``up'' quarks and one ``down'' quark (charges $+\tfrac23e,+\tfrac23e,-\tfrac13e$, total $+e$), the neutron
from one up and two down ($+\tfrac23e,-\tfrac13e,-\tfrac13e$, total $0$). In $\beta^-$ decay a single
down quark turns into an up quark through the weak force, emitting an electron and an antineutrino:
$d\to u+e^-+\bar\nu_e$. Charge is conserved ($-\tfrac13=+\tfrac23-1+0$), and so is lepton number
($0=+1-1$). In $\beta^+$ decay an up quark turns into a down quark. None of this is needed for the
calculations in this chapter, but it explains why the weak force, and only the weak force, can change a
neutron into a proton.
\end{remarkbox}

\section{Three Kinds of Radiation}\label{sec:decay-three}

With the particles named, the radiation itself can be sorted. Rutherford separated the emissions of
uranium and radium by two tests: how far they penetrate matter, and how they bend in electric and
magnetic fields (Figures~\ref{fig:decay-penetration} and~\ref{fig:decay-deflection}). He named the three
kinds after the first three letters of the Greek alphabet. Table~\ref{tab:decay-radiation} collects what
is now known about each.

\begin{tablebox}{Alpha, beta and gamma radiation}
\centering
\renewcommand{\arraystretch}{1.15}
\begin{tabular}{@{}>{\raggedright}p{2.3cm}>{\raggedright}p{3.6cm}>{\raggedright}p{3.6cm}>{\raggedright\arraybackslash}p{3.6cm}@{}}
\toprule
 & \textbf{alpha} ($\alpha$) & \textbf{beta} ($\beta^\mp$) & \textbf{gamma} ($\gamma$)\\
\midrule
Nature & helium-4 nucleus $\nuc42{He}$ & electron $e^-$ (or positron $e^+$) & high-energy photon\\
Charge & $+2e$ & $-e$ (or $+e$) & $0$\\
Mass & $\SI{4.0015}{\amu}$ & $\SI{0.000549}{\amu}$ & $0$\\
Typical energy & $4$--$\SI{9}{\MeV}$, \emph{discrete} & up to a few MeV, \emph{continuous} & keV to a few MeV, \emph{discrete}\\
Typical speed & about $0.05c$ & up to $0.99c$ & $c$\\
Ionising power & very high & moderate & low\\
Range in air & a few cm & up to metres & very long (weakened, not stopped)\\
Stopped by & a sheet of paper, dead skin & a few mm of aluminium & several cm of lead reduce it greatly\\
\bottomrule
\end{tabular}\label{tab:decay-radiation}
\end{tablebox}

The speeds follow from the energies. A $\SI{5}{\MeV}$ $\alpha$ particle has kinetic energy only
$5/3727=0.13\%$ of its rest energy, so the non-relativistic $K=\tfrac12mv^2$ of \chref{ch:energy} applies
and $v/c=\sqrt{2K/mc^2}=\sqrt{2(5)/3727}=0.052$. An electron of a few MeV has a kinetic energy several
times its rest energy of $\SI{0.511}{\MeV}$, so it moves at nearly $c$, and its kinematics would need
relativity; this chapter avoids electron kinematics except where only energy conservation is used.

\begin{figure}[htbp]
\centering
\tikzfig{decay_penetration}
\figcaption{fig:decay-penetration}{Penetrating power. Heavy, doubly charged $\alpha$ particles ionise
densely and lose their energy within a few centimetres of air or a sheet of paper. Electrons penetrate
farther, to a few millimetres of aluminium. Uncharged photons interact only occasionally, so their
intensity falls exponentially with the thickness of lead and never reaches exactly zero: a weakened beam
always gets through.}
\end{figure}

\begin{keybox}{Ionisation and range}
An ionising particle loses energy by pulling electrons off the atoms it passes. The impulse given to each
electron grows with the particle's charge and with the \emph{time} it spends nearby. A slow, doubly charged
$\alpha$ particle lingers and pulls hard, creating tens of thousands of ion pairs per centimetre of air; a
fast electron carries a single charge and races past. Dense ionisation means rapid energy loss and therefore
a short range:
\begin{center}
heavier, slower, more highly charged $\;\Longrightarrow\;$ denser ionisation $\;\Longrightarrow\;$ shorter range.
\end{center}
An $\alpha$ source is therefore harmless outside the body (the dead outer layer of skin stops it) but very
dangerous if inhaled or swallowed, where it deposits all its energy in a few hundredths of a millimetre of
living tissue.
\end{keybox}

\begin{trapbox}{gamma rays are weakened, not stopped}
``Lead stops gamma rays'' is a convenient slogan and a wrong statement. A photon either passes a layer
untouched or is removed by a single interaction, so each centimetre of lead removes the same
\emph{fraction} of the beam. The intensity falls exponentially and never reaches zero; shielding is
specified by how many times it divides the intensity, not by a thickness that blocks it. The
mathematics of such exponential falls is developed in \chref{ch:decaylaw}.
\end{trapbox}

\textbf{Bending in a magnetic field.} The second of Rutherford's tests uses the magnetic force on a moving
charge. A charge $q$ moving with velocity $\vv$ through a magnetic field $\B$ feels the force
$\vF=q\,\vv\times\B$, perpendicular to both (the cross product is defined in \chref{ch:toolkit},
Section~\ref{sec:vecalg}). This law belongs to electromagnetism and is stated here without proof. Because
the force is always perpendicular to $\vv$, it does no work and changes only the direction of motion. When
$\vv$ is perpendicular to $\B$ the particle therefore moves on a circle, and the magnetic force of size
$|q|vB$ supplies the centripetal force $mv^2/r$ of \chref{ch:circdyn}:
\begin{equation}\label{eq:decay-radius}
|q|vB=\frac{mv^2}{r}\quad\Longrightarrow\quad \boxed{\;r=\frac{mv}{|q|B}=\frac{p}{|q|B},\;}
\end{equation}
where $p$ is the magnitude of the momentum. (For fast electrons the result $r=p/|q|B$ still holds,
with the relativistic momentum; it is quoted, not derived, here.) A large momentum and a small charge
make a large, gently curved circle. Positive and negative charges bend in opposite directions, and
photons, with no charge, do not bend at all (Figure~\ref{fig:decay-deflection}).

\begin{figure}[htbp]
\centering
\tikzfig{decay_deflection}
\figcaption{fig:decay-deflection}{Deflection in a magnetic field directed into the page. The force
$q\vv\times\B$ bends positive $\alpha$ particles and negative electrons in opposite directions, and leaves
uncharged photons undeflected. The radius $r=p/|q|B$ is far larger for the $\alpha$ particle, because its
momentum is much larger: at $\SI{1}{T}$ a $\SI{5.3}{\MeV}$ $\alpha$ circles with $r=\SI{33}{cm}$
(Example~\ref{ex:decay-alpharadius}) and a $\SI{1}{\MeV}$ electron with $r\approx\SI{5}{mm}$. The drawing
exaggerates the $\alpha$ bend about tenfold.}
\end{figure}

\begin{exbox}{An alpha particle in a magnetic field}\label{ex:decay-alpharadius}
Find the radius of the path of a $\SI{5.30}{\MeV}$ $\alpha$ particle (mass $\SI{6.645e-27}{kg}$, charge
$+2e$) moving perpendicular to a magnetic field of $\SI{1.00}{T}$, and its speed as a fraction of $c$.

\stepn{1} \textbf{Momentum.} The $\alpha$ is non-relativistic ($K\ll mc^2$), so $p=\sqrt{2mK}$ with
$K=5.30\times\SI{1.602e-13}{J}$:
$p=\sqrt{2(6.645\times10^{-27})(8.49\times10^{-13})}=\SI{1.06e-19}{\kilo\gram\metre\per\second}$.
\stepn{2} \textbf{Radius.} From Eq.~\eqref{eq:decay-radius}, $r=p/(2eB)=
\dfrac{1.06\times10^{-19}}{2(1.602\times10^{-19})(1.00)}=\SI{0.332}{m}$.
\stepn{3} \textbf{Speed.} $v=p/m=\SI{1.60e7}{\metre\per\second}=0.053c$, matching the estimate above.
\ans{$r\approx\SI{33}{cm}$, $v=0.053c$}
\end{exbox}

Every radiation detector exploits the ionisation described above, and one of them gave radioactivity its
familiar sound.

\begin{remarkbox}{Seeing the invisible}
A \emph{Geiger--M\"uller tube} (Figure~\ref{fig:decay-geiger}) is a gas-filled cylinder with a thin wire
at high voltage along its axis. A single ionising particle frees a few electrons, which accelerate towards
the wire and knock out more electrons in an avalanche; each particle produces one sharp pulse, heard as a
click. A \emph{cloud chamber} contains supersaturated vapour that condenses on ions, so a particle leaves
a visible trail. The trails mirror Table~\ref{tab:decay-radiation}: $\alpha$ particles make short, thick,
straight tracks (dense ionisation, heavy, hard to deflect), electrons make long, thin, wandering ones.
Modern instruments use scintillator crystals, which flash when struck, and semiconductor diodes, which
measure the energy deposited.
\end{remarkbox}

\begin{figure}[htbp]
\centering
\tikzfig{decay_geiger}
\figcaption{fig:decay-geiger}{A Geiger--M\"uller tube. One ionising particle frees a few electrons in the
gas; the strong field near the thin central wire accelerates them so violently that each frees more, and
the avalanche gives a pulse large enough to count. The pulse is the same whatever the particle's energy,
so a Geiger counter counts particles but cannot measure their energies.}
\end{figure}

\section{The Rules Every Decay Obeys}\label{sec:decay-rules}

Knowing what comes out is not yet enough to predict a decay. One needs the rules that every decay,
whatever its kind, must satisfy. A nuclear decay is a reaction
$\text{parent}\to\text{daughter}+\text{emitted particles}$, and it must conserve:
\begin{enumerate}[label=\textbf{\arabic*.}]
\item \textbf{electric charge}, so the lower numbers (charges, counting $e^-$ as $-1$ and $e^+$ as $+1$)
      balance;
\item \textbf{nucleon number}, so the upper numbers balance;
\item \textbf{energy}, including rest-mass energy, and \textbf{linear momentum}, exactly as in the
      collisions and explosions of \chref{ch:collisions};
\item \textbf{lepton number} (Section~\ref{sec:decay-particles}): an emitted electron must be accompanied
      by an antineutrino, and an emitted positron or a captured electron by a neutrino.
\end{enumerate}
The energy rule needs a number attached to it, and that number is the $Q$-value.

\begin{defbox}{$Q$-value}
The $Q$-value of a decay is the rest-mass energy that disappears:
\[
Q=\big(m_{\text{initial}}-m_{\text{final}}\big)c^2 .
\]
It appears as kinetic energy shared among the products. A decay can occur spontaneously only if $Q>0$.
\end{defbox}

The reason is energy conservation itself. The parent is at rest, with energy $m_{\text{initial}}c^2$. The
products have energy $m_{\text{final}}c^2$ plus their kinetic energies, which cannot be negative. So the
products can be formed only if $m_{\text{final}}<m_{\text{initial}}$, and the difference, $Q$, is exactly
the kinetic energy they share. A decay with $Q>0$ is \emph{allowed}; whether it happens quickly or in
billions of years is a separate question.

\why{Mass tables list the masses $M$ of neutral \emph{atoms}, electrons included, not the masses $m$ of bare
nuclei. Each decay below is first written with nuclear masses, where the physics is plain, and then
converted to atomic masses, where the numbers are. The conversion is $M=m+Zm_e$, neglecting the binding
energies of the electrons, whose difference between parent and daughter is at most tens of keV, far below
the MeV $Q$-values. Whether the electron masses cancel or not is the commonest source of wrong answers in
this chapter.}

\section{Alpha Decay}\label{sec:decay-alpha}

The first kind of radiation to examine in detail is the one that changes the nucleus most. In
$\alpha$ decay the nucleus ejects a helium-4 nucleus:
\[
\nuc{A}{Z}{X}\;\longrightarrow\;\nuc{A-4}{Z-2}{Y}+\nuc42{He},\qquad\text{for example}\quad
\nuc{226}{88}{Ra}\to\nuc{222}{86}{Rn}+\nuc42{He}.
\]
Here X is the parent and Y the daughter, $A$ the mass number and $Z$ the atomic number. The daughter sits
two places lower in the periodic table, and its mass number is four less.

\begin{derivbox}{The $Q$-value of $\alpha$ decay from atomic masses}
With nuclear masses, $Q=[m(\mathrm{X})-m(\mathrm{Y})-m(\nuc42{He})]c^2$. Add $Zm_e$ to the first term, and
$(Z-2)m_e+2m_e=Zm_e$ to the last two together. The added electrons balance exactly, so atomic masses may be
used directly:
\begin{equation}\label{eq:decay-Qalpha}
\boxed{\;Q_\alpha=\big[M(\mathrm{X})-M(\mathrm{Y})-M(\iso4{He})\big]c^2 .\;}
\end{equation}
\end{derivbox}

How is that energy shared between the two products? This is an explosion of the kind treated in
\chref{ch:collisions}, a perfectly inelastic collision run backwards: no external force acts during the
decay, so the total momentum, zero before, is zero after.

\begin{derivbox}{How the energy is shared}
The parent is at rest, so the $\alpha$ particle and the daughter fly apart with equal and opposite
momenta, $\vp_\alpha=-\vp_{\mathrm{Y}}$, of common magnitude $p$. Both are slow compared with light (the
$\alpha$ carries about $0.1\%$ of its rest energy), so $K=p^2/2m$ for each, and
\[
\frac{K_\alpha}{K_{\mathrm{Y}}}=\frac{p^2/2m_\alpha}{p^2/2m_{\mathrm{Y}}}=\frac{m_{\mathrm{Y}}}{m_\alpha}\approx\frac{A-4}{4},
\qquad K_\alpha+K_{\mathrm{Y}}=Q,
\]
where the masses have been replaced by their mass numbers, accurate to about $0.1\%$. Solving the two
equations:
\begin{equation}\label{eq:decay-Kalpha}
\boxed{\;K_\alpha=Q\,\frac{A-4}{A},\qquad K_{\mathrm{Y}}=Q\,\frac{4}{A}.\;}
\end{equation}
The lighter particle takes almost all the energy, just as the bullet, not the rifle, carries most of the
kinetic energy when a gun is fired. Because a two-body decay fixes the share uniquely, every $\alpha$ from
a given transition has the \emph{same} energy: the $\alpha$ spectrum is a sharp line
(Figure~\ref{fig:decay-betaspec}, left).
\end{derivbox}

\begin{exbox}{Radium-226}\label{ex:decay-ra226}
Find the energy released in $\nuc{226}{88}{Ra}\to\nuc{222}{86}{Rn}+\nuc42{He}$ and its share between the two
products, given $M(\iso{226}{Ra})=\SI{226.025410}{\amu}$, $M(\iso{222}{Rn})=\SI{222.017578}{\amu}$ and
$M(\iso4{He})=\SI{4.002603}{\amu}$.

\stepn{1} \textbf{$Q$-value.} By Eq.~\eqref{eq:decay-Qalpha}, the mass lost is
$226.025410-222.017578-4.002603=\SI{0.005229}{\amu}$, so $Q=(0.005229)(931.494)=\SI{4.871}{\MeV}$.
\stepn{2} \textbf{Sharing.} By Eq.~\eqref{eq:decay-Kalpha}, $K_\alpha=4.871\times\tfrac{222}{226}=\SI{4.784}{\MeV}$
and $K_{\mathrm{Rn}}=4.871\times\tfrac{4}{226}=\SI{0.086}{\MeV}$.
\stepn{3} \textbf{Check.} $4.784+0.086=4.870$, equal to $Q$ within rounding; the radon recoils with less
than $2\%$ of the energy. \ans{$Q=\SI{4.871}{\MeV}$; $\SI{4.784}{\MeV}$ and $\SI{0.086}{\MeV}$}
\end{exbox}

\begin{exbox}{Polonium-210}\label{ex:decay-po210}
Given $M(\iso{210}{Po})=\SI{209.982874}{\amu}$, $M(\iso{206}{Pb})=\SI{205.974465}{\amu}$ and
$M(\iso4{He})=\SI{4.002603}{\amu}$, find $Q$ for the $\alpha$ decay of polonium-210 and the kinetic energies
of both products.

The mass lost is $\Delta M=\SI{0.005806}{\amu}$, so $Q=(0.005806)(931.494)=\SI{5.408}{\MeV}$. Then
$K_\alpha=5.408\times\tfrac{206}{210}=\SI{5.305}{\MeV}$, matching the measured line at \SI{5.304}{\MeV} in
Figure~\ref{fig:decay-betaspec}, and $K_{\mathrm{Pb}}=5.408\times\tfrac{4}{210}=\SI{0.103}{\MeV}$.
\ans{$Q=\SI{5.408}{\MeV}$; $\SI{5.305}{\MeV}$ and $\SI{0.103}{\MeV}$}
\end{exbox}

\begin{trapbox}{the $\alpha$ particle does not carry all of $Q$}
It is tempting to equate the measured $\alpha$ energy with the $Q$-value. The daughter's recoil takes
$4/A$ of $Q$, about $\SI{0.1}{\MeV}$ for heavy nuclei: small, but larger than the precision of an $\alpha$
energy measurement, and in a solid source large enough to knock the recoiling atom out of its place.
Going backwards from a measured $K_\alpha$, $Q=K_\alpha\,A/(A-4)$.
\end{trapbox}

The formulas say how much energy is released; they do not say why the nucleus emits this particular
cluster, nor how it gets out.

\begin{remarkbox}{Why an $\alpha$, and how it escapes}
Why should a nucleus emit a bound cluster of four nucleons rather than single protons or neutrons? Because
$\iso4{He}$ is so tightly bound ($\Eb=\SI{28.3}{\MeV}$, \chref{ch:nucleus}, Section~\ref{sec:nuc-binding}) that ejecting it is often
energetically allowed ($Q>0$) when ejecting a single nucleon is not ($Q<0$). Escaping is still puzzling.
For an $\alpha$ particle touching a thorium-234 nucleus, the centres are about
$1.2(234^{1/3}+4^{1/3})=\SI{9.3}{\fm}$ apart, and the electric potential energy there is
$(2)(90)(\SI{1.44}{\MeV.\fm})/(\SI{9.3}{\fm})\approx\SI{28}{\MeV}$; yet the $\alpha$ from uranium-238
emerges with only $\SI{4.2}{\MeV}$. Classically it could never have climbed over this \emph{Coulomb
barrier}. In 1928 George Gamow explained that the $\alpha$ is a quantum wave and \emph{tunnels} through
the barrier: because a quantum particle is spread out like a wave, there is a small but non-zero
probability of finding it on the far side of a barrier it could never climb classically (stated without
proof; it needs quantum mechanics). The tunnelling probability is exponentially sensitive to energy. That
explains the Geiger--Nuttall rule: as the $\alpha$ energy rises from $\SI{4.2}{\MeV}$ (uranium-238) to
$\SI{7.7}{\MeV}$ (polonium-214), the half-life falls from $4.5\times10^{9}$ years to $\SI{164}{\micro\second}$, a
factor of about $10^{21}$.
\end{remarkbox}

\section{Beta Decay and Its Relatives}\label{sec:decay-beta}

Alpha decay removes nucleons; beta decay rearranges them. It changes a neutron into a proton, or the
reverse, inside the nucleus. The mass number is unchanged, so parent and daughter are \emph{isobars}
(same $A$, \chref{ch:nucleus}). There are three varieties, pictured in Figure~\ref{fig:decay-betaproc}.

\textbf{Beta-minus ($\beta^-$).} A nucleus with too many neutrons converts one of them:
\[
n\to p+e^-+\bar\nu_e,\qquad \nuc{A}{Z}{X}\to\nuc{A}{Z+1}{Y}+e^-+\bar\nu_e,\qquad\text{for example}\quad
\nuc{14}{6}{C}\to\nuc{14}{7}{N}+e^-+\bar\nu_e .
\]
\textbf{Beta-plus ($\beta^+$, positron emission).} A nucleus with too many protons converts one of them:
\[
p\to n+e^++\nu_e,\qquad \nuc{A}{Z}{X}\to\nuc{A}{Z-1}{Y}+e^++\nu_e,\qquad\text{for example}\quad
\nuc{22}{11}{Na}\to\nuc{22}{10}{Ne}+e^++\nu_e .
\]
\textbf{Electron capture (EC).} The nucleus absorbs one of its own inner (K-shell) electrons:
\[
p+e^-\to n+\nu_e,\qquad \nuc{A}{Z}{X}+e^-\to\nuc{A}{Z-1}{Y}+\nu_e,\qquad\text{for example}\quad
\nuc74{Be}+e^-\to\nuc73{Li}+\nu_e .
\]
The vacancy left in the K shell is filled by an outer electron, which emits a characteristic X-ray. Since
the neutrino escapes unseen, that X-ray is how electron capture is detected.

\begin{figure}[htbp]
\centering
\tikzfig{decay_betaproc}
\figcaption{fig:decay-betaproc}{The three beta processes, seen inside the nucleus. In each, one nucleon
changes identity and the mass number is unchanged. Charge and lepton number balance in every case:
$\beta^-$ creates an electron ($L_e=+1$) and an antineutrino ($L_e=-1$); $\beta^+$ creates a positron ($-1$)
and a neutrino ($+1$); electron capture destroys an electron ($+1$) and creates a neutrino ($+1$), and the
refilling of the K shell gives an X-ray.}
\end{figure}

The electrons and positrons are \emph{created} at the moment of decay, much as a photon is created when an
atom de-excites; they were never inside the nucleus. A nucleus could not hold an electron anyway. By
Heisenberg's uncertainty principle (stated without proof), a particle confined to a region of size
$\Delta x$ must have a spread of momentum of at least about $\hbar/\Delta x$, where $\hbar=h/2\pi$.
Confining an electron within $\Delta x\approx\SI{10}{\fm}$ therefore gives it momentum of order
$\hbar/\Delta x$, that is, $pc\approx\hbar c/\Delta x=(\SI{197}{\MeV.\fm})/(\SI{10}{\fm})\approx\SI{20}{\MeV}$.
That is far more than the energy of any $\beta$ electron ever observed (at most a few MeV), and far more
than the electric attraction of the nucleus could hold.

The force responsible is the \emph{weak nuclear force}, the only interaction that can change a neutron
into a proton or the reverse. It is very feeble and very short-ranged, which is why $\beta$ decays are
slow: typical $\beta$ half-lives are seconds to years, whereas a nucleus rearranging under the strong
force takes about $\SI{e-22}{s}$.

Writing the $Q$-values for the three processes is where the warning about atomic and nuclear masses
(Section~\ref{sec:decay-rules}) matters most.

\begin{derivbox}{$Q$-values with atomic masses: the electron trap}
Let $m$ denote nuclear masses and $M$ atomic masses, so $M(\mathrm{X})=m(\mathrm{X})+Zm_e$ (electron
binding energies neglected).
\begin{itemize}
\item $\beta^-$: $Q=[m(\mathrm{X})-m(\mathrm{Y})-m_e]c^2=\big[(M(\mathrm{X})-Zm_e)-(M(\mathrm{Y})-(Z+1)m_e)-m_e\big]c^2$.
      The electrons cancel:
      \begin{equation}\label{eq:decay-Qbminus}
      \boxed{\;Q_{\beta^-}=[M(\mathrm{X})-M(\mathrm{Y})]c^2.\;}
      \end{equation}
      The daughter atom needs $Z+1$ electrons; the parent's $Z$ plus the emitted one make up the count.
\item $\beta^+$: $Q=[m(\mathrm{X})-m(\mathrm{Y})-m_e]c^2=\big[(M(\mathrm{X})-Zm_e)-(M(\mathrm{Y})-(Z-1)m_e)-m_e\big]c^2$,
      which leaves two electron masses:
      \begin{equation}\label{eq:decay-Qbplus}
      \boxed{\;Q_{\beta^+}=[M(\mathrm{X})-M(\mathrm{Y})-2m_e]c^2.\;}
      \end{equation}
      One $m_e$ is the positron created; the other is the parent's surplus atomic electron, which the
      daughter atom no longer needs.
\item EC: $Q=[m(\mathrm{X})+m_e-m(\mathrm{Y})]c^2=\big[(M(\mathrm{X})-Zm_e)+m_e-(M(\mathrm{Y})-(Z-1)m_e)\big]c^2$:
      \begin{equation}\label{eq:decay-Qec}
      \boxed{\;Q_{\mathrm{EC}}=[M(\mathrm{X})-M(\mathrm{Y})]c^2.\;}
      \end{equation}
      Strictly, the daughter atom is left with a K-shell vacancy; the neutrino carries $Q_{\mathrm{EC}}$ minus
      the K-shell binding energy, and that small remainder (keV) leaves as the X-ray.
\end{itemize}
The neutrino's mass, below $\SI{1}{\eV}/c^2$, is neglected throughout.
\end{derivbox}

\begin{keybox}{The $2m_ec^2$ threshold}
Positron emission requires the parent \emph{atom} to outweigh the daughter atom by more than
$2m_ec^2=\SI{1.022}{\MeV}$. Electron capture requires only that the parent atom be heavier at all. A
proton-rich nucleus whose atomic mass exceeds its daughter's by less than $\SI{1.022}{\MeV}/c^2$ can
therefore decay \emph{only} by electron capture. Beryllium-7 is the classic example
(Practice Problem~\ref{prob:decay-be7}).
\end{keybox}

\begin{trapbox}{the same mass difference gives two different $Q$-values}
For a given pair of atoms, $\beta^-$ and EC use the bare atomic mass difference, while $\beta^+$ subtracts
$\SI{1.022}{\MeV}$. A common error subtracts nothing (treating $\beta^+$ like $\beta^-$), or only
$\SI{0.511}{\MeV}$ (counting the positron but forgetting the surplus atomic electron). Both errors
disappear if the masses are first written as nuclear masses.
\end{trapbox}

\begin{exbox}{Sodium-22: two channels}\label{ex:decay-na22}
The atomic masses are $M(\iso{22}{Na})=\SI{21.994437}{\amu}$ and $M(\iso{22}{Ne})=\SI{21.991385}{\amu}$. Find
$Q$ for $\beta^+$ decay and for electron capture.

The atomic mass difference is $\SI{0.003052}{\amu}$, which is $(0.003052)(931.494)=\SI{2.843}{\MeV}$. By
Eq.~\eqref{eq:decay-Qec}, $Q_{\mathrm{EC}}=\SI{2.843}{\MeV}$; by Eq.~\eqref{eq:decay-Qbplus},
$Q_{\beta^+}=2.843-1.022=\SI{1.821}{\MeV}$. Both channels are open, and sodium-22 uses both.
\ans{$Q_{\beta^+}=\SI{1.821}{\MeV}$, $Q_{\mathrm{EC}}=\SI{2.843}{\MeV}$}
\end{exbox}

\begin{exbox}{Checking a decay against every rule}\label{ex:decay-na22check}
Is $\nuc{22}{11}{Na}\to\nuc{22}{10}{Ne}+e^++\nu_e$ allowed?

Go through the rules of Section~\ref{sec:decay-rules} one at a time. Charge: $11=10+1+0$. Nucleons:
$22=22+0+0$. Lepton number: $0=0-1+1$. Energy: $Q_{\beta^+}=\SI{1.821}{\MeV}>0$
(Example~\ref{ex:decay-na22}). Momentum: three products can share zero total momentum in many ways. All
five pass, so the decay is allowed, and it is observed. \ans{allowed}
\end{exbox}

\begin{remarkbox}{The free neutron}
Outside a nucleus a neutron is unstable. Its mass exceeds that of a hydrogen atom, so
$n\to p+e^-+\bar\nu_e$ releases $Q=(1.008665-1.007825)(931.494)=\SI{0.782}{\MeV}$ (a $\beta^-$ decay, so
the hydrogen \emph{atomic} mass is the right one to use). Its mean life is about $\SI{880}{s}$, roughly
fifteen minutes. The reverse, $p\to n+e^++\nu_e$, is impossible for a free proton, which is lighter.
Inside a nucleus it can happen, because the difference in binding energy between parent and daughter
supplies the missing energy.
\end{remarkbox}

\section{The Missing Energy and the Neutrino}\label{sec:decay-neutrino}

The neutrino has appeared in every $\beta$ equation so far, yet no detector registers it. Why believe in
it? The evidence is the energy of the electrons. If $\beta^-$ decay produced only two bodies, daughter plus
electron, momentum conservation would fix the electron's energy exactly, as it fixes the $\alpha$ energy
(Eq.~\eqref{eq:decay-Kalpha}). Instead, experiments from 1914 to the late 1920s found a \emph{continuous}
spread of electron energies, from zero up to a maximum equal to $Q$ (Figure~\ref{fig:decay-betaspec}). On
average about two-thirds of the energy seemed to vanish. Niels Bohr was prepared to give up energy
conservation in individual nuclear processes. In 1930 Wolfgang Pauli proposed what he called a desperate
remedy: a third, neutral, nearly massless particle carries off the missing energy and momentum. Fermi named
it the \emph{neutrino} (``little neutral one'') and built the theory of beta decay around it. It interacts
so weakly that a typical MeV neutrino could pass through a light-year of lead. It was finally detected in
1956 by Cowan and Reines beside a nuclear reactor.

\begin{figure}[htbp]
\centering
\tikzfig{decay_betaspec}
\figcaption{fig:decay-betaspec}{Two-body versus three-body decay. Left: every $\alpha$ from polonium-210
carries $\SI{5.30}{\MeV}$ (Example~\ref{ex:decay-po210}). Right: the electron energy in $\beta^-$ decay is
spread continuously from $0$ to $Q$, and the antineutrino carries the rest. The curve is the simplest
theoretical shape for $Q=\SI{1}{\MeV}$, ignoring the pull of the nucleus on the outgoing electron; its mean
is $0.40Q$. That pull slows real $\beta^-$ electrons, and their mean is nearer $Q/3$.}
\end{figure}

\begin{keybox}{Three bodies share energy freely}
With three products, the momenta need only add to zero, that is, form a closed triangle, and infinitely
many triangles are possible (Figure~\ref{fig:decay-momentum}). The share of energy therefore varies from
decay to decay. Neglecting the tiny recoil of the heavy daughter,
\[
K_e+E_{\bar\nu}=Q,\qquad 0\le K_e\le Q,
\]
where $K_e$ is the electron's kinetic energy and $E_{\bar\nu}$ the antineutrino's energy. The maximum
electron energy, the \textbf{endpoint}, equals $Q$ and occurs when the antineutrino is left with
essentially nothing.
\end{keybox}

\begin{exbox}{Sharing in a carbon-14 decay}\label{ex:decay-c14share}
In a carbon-14 decay ($Q=\SI{0.156}{\MeV}$) an electron is detected with $K_e=\SI{0.100}{\MeV}$. What
energy does the antineutrino carry?

Neglecting the nuclear recoil (a few eV, as shown next), $E_{\bar\nu}=Q-K_e=0.156-0.100=\SI{0.056}{\MeV}$.
\ans{$\SI{0.056}{\MeV}$}
\end{exbox}

The neglected recoil really is tiny. Even an electron at the carbon-14 endpoint ($\SI{0.156}{\MeV}$) has
momentum $pc=\sqrt{K^2+2Km_ec^2}=\SI{0.43}{\MeV}$ (the relativistic relation, quoted here), and a
nitrogen-14 nucleus taking that momentum moves off with $(pc)^2/2Mc^2\approx\SI{7}{\eV}$, a
twenty-thousandth of $Q$. The daughter takes up momentum but almost no energy.

\begin{figure}[htbp]
\centering
\tikzfig{decay_momentum}
\figcaption{fig:decay-momentum}{Why two-body and three-body decays differ. All the arrows are momenta.
Left: in $\alpha$ decay the two momenta must be equal and opposite, which fixes the energy of each. Right:
in $\beta^-$ decay any three momenta that add to zero are allowed (both examples close exactly), so the
electron's share varies continuously from decay to decay.}
\end{figure}


\begin{trapbox}{endpoint versus typical energy}
Tables quote the $\beta$ endpoint $Q$, but almost no electron has that energy; the average is roughly
$Q/3$. When computing the heat deposited by a $\beta$ source, use the average electron energy: the
neutrinos' share escapes (Practice Problem~\ref{prob:decay-mean}).
\end{trapbox}

\section{Gamma Decay}\label{sec:decay-gamma}

Alpha and beta decays change the nucleus into a different nuclide; they often leave it with energy to
spare. The daughter is then in an \emph{excited} nuclear state and, like an excited atom, drops to lower
states by emitting photons. Nuclear energy levels are spaced by keV to MeV rather than eV, so the photons
are gamma rays:
\[
\nuc{A}{Z}{X}^{*}\to\nuc{A}{Z}{X}+\gamma ,
\]
where the star marks the excited state. Neither $Z$ nor $A$ changes: gamma decay does not turn one element
into another. Because the levels are discrete, gamma energies are sharp lines that fingerprint the
nuclide. Occasionally the excited nucleus hands its energy directly to one of its own inner atomic
electrons, which is ejected with a definite energy instead of a photon being emitted. This competing
process, \textbf{internal conversion}, is followed by X-rays as the vacancy is filled.

Momentum conservation applies here too: the photon carries momentum $E_\gamma/c$, so the nucleus recoils
with the same momentum and kinetic energy $(E_\gamma/c)^2/2M=E_\gamma^2/2Mc^2$. For a $\SI{1.333}{\MeV}$
photon from a nickel-60 nucleus ($Mc^2=60\times\SI{931.5}{\MeV}$) this is
$1.333^2/(2\times55\,890)\,\si{\MeV}=\SI{16}{\eV}$, so to excellent accuracy the gamma energy equals the
spacing of the levels.

\begin{figure}[htbp]
\centering
\tikzfig{decay_co60scheme}
\figcaption{fig:decay-co60}{Decay scheme of cobalt-60, the workhorse of radiotherapy. Energy runs
upwards. The $\beta^-$ decay lands on an excited state of nickel-60, which cascades to the ground state by
emitting two gamma rays. The energy bookkeeping closes:
$0.316+1.173+1.333=\SI{2.822}{\MeV}=Q$.}
\end{figure}

Figure~\ref{fig:decay-co60} shows how $\beta$ and $\gamma$ decay combine in one of the most used
radioactive sources. The $\beta^-$ decay does not go to the ground state of nickel-60 but to a level
$\SI{2.506}{\MeV}$ above it, so the electrons' endpoint is not $Q$ but $Q-\SI{2.506}{\MeV}$
(Practice Problem~\ref{prob:decay-co60}).

\begin{remarkbox}{Technetium-99m, a nucleus that waits}
Most excited nuclear states decay in under $\SI{e-12}{s}$. A few, called \emph{isomers} (marked m), are
trapped for much longer, because the transition requires a large change of angular momentum. Nuclei,
like electrons, carry a quantised intrinsic angular momentum, or \emph{spin}, and a photon finds a large
change of it very hard to carry away. Technetium-99m sits in its excited state with a half-life of
$\SI{6.0}{h}$ before emitting a single $\SI{140}{\keV}$ gamma ray. That is long enough to prepare and inject
a radiopharmaceutical, short enough that the patient's dose fades within a day, and the gamma energy
passes easily out of the body to a camera. It is used in tens of millions of diagnostic scans every
year (\chref{ch:nucapps}).
\end{remarkbox}

\section{Which Nuclei Are Unstable, and How They Decay}\label{sec:decay-stability}

The decay modes are now known; the remaining question is which nuclei use which. The answer is a map.
Figure~\ref{fig:decay-nzchart} plots the naturally occurring nuclides by proton number $Z$ and neutron
number $A-Z$. Stable nuclei form a narrow band, the \textbf{valley of stability}. It follows $A-Z=Z$ for
light nuclei and bends towards neutron excess for heavy ones, because extra neutrons add nuclear
attraction without adding electric repulsion (\chref{ch:nucleus}).

\begin{figure}[htbp]
\centering
\tikzfig{decay_nzchart}
\figcaption{fig:decay-nzchart}{The naturally occurring nuclides, plotted by proton number $Z$ and neutron
number $A-Z$. The stable ones (filled dots) form the valley of stability, which bends away from the line
$A-Z=Z$ towards neutron excess as $Z$ grows. The open circles are radioactive nuclides so long-lived that
they have survived since the Earth formed, such as potassium-40, thorium-232 and uranium-238. Nuclei above
the valley have too many neutrons and undergo $\beta^-$ decay; nuclei below it have too many protons and
undergo $\beta^+$ decay or electron capture. Lead ($Z=82$) is the heaviest element with stable nuclides;
heavier nuclei shed mass mainly by $\alpha$ decay. (Bismuth-209, long counted as stable, was found in 2003
to $\alpha$-decay with a half-life of about $2\times10^{19}$ years.)}
\end{figure}

\textbf{Why $\beta$ decay always heads for the valley.} Fix the mass number $A$ and compare all its
isobars, the nuclei that differ only in how the $A$ nucleons are split between protons and neutrons. A
$\beta$ decay turns one isobar into its neighbour, and by Eqs.~\eqref{eq:decay-Qbminus}--\eqref{eq:decay-Qec}
it can happen only if the neighbour's atomic mass is lower. So the question is how the atomic mass varies
with $Z$ at fixed $A$, and the liquid-drop formula of \chref{ch:nucleus} (Figure~\ref{fig:nuc-semfterms}) answers it.

\begin{derivbox}{Why the isobar masses lie on a parabola}
Build the atom from $Z$ hydrogen atoms (each of atomic mass $M_{\mathrm{H}}=M(\iso{1}{H})$) and $A-Z$ neutrons, and subtract the binding energy:
$M(Z)c^2=ZM_{\mathrm{H}}c^2+(A-Z)m_nc^2-\Eb$. In the liquid-drop formula only two terms of $\Eb$ depend
on $Z$ at fixed $A$: the Coulomb term $-a_CZ(Z-1)/A^{1/3}$ and the asymmetry term $-a_A(A-2Z)^2/A$, with
$a_C\approx\SI{0.71}{\MeV}$ and $a_A\approx\SI{23.7}{\MeV}$ (pairing aside). Both are quadratic in $Z$, so
\[
M(Z)c^2=\text{const}-(m_n-M_{\mathrm{H}})c^2\,Z+\frac{a_CZ(Z-1)}{A^{1/3}}+\frac{a_A(A-2Z)^2}{A},
\]
a parabola in $Z$ that opens upwards. Its minimum, where $dM/dZ=0$, is at
\[
Z_0=\frac{4a_A+(m_n-M_{\mathrm{H}})c^2+a_CA^{-1/3}}{8a_A/A+2a_CA^{-1/3}} .
\]
Without the Coulomb term this would be $Z_0\approx A/2$ (the asymmetry term alone wants equal numbers);
the Coulomb term, growing with $Z^2$, pushes the minimum to fewer protons, more so for large $A$. For
$A=101$, with $(m_n-M_{\mathrm{H}})c^2=\SI{0.782}{\MeV}$, the formula gives $Z_0=43.9$; the measured lightest
isobar is ruthenium-101, $Z=44$.
\end{derivbox}

Every isobar on the neutron-rich side of the minimum can roll one step downhill by $\beta^-$ decay, every
isobar on the proton-rich side by $\beta^+$ decay or electron capture, and only the isobar at the bottom is
stable (Figure~\ref{fig:decay-isobar}). Plotting the bottoms of the parabolas for every $A$ traces out the
valley of Figure~\ref{fig:decay-nzchart}. For even $A$ the pairing effect of \chref{ch:nucleus} splits the
curve into two parallel parabolas, one for even-$Z$ and one for odd-$Z$ nuclei, and sometimes two isobars
are both stable.

\begin{figure}[htbp]
\centering
\tikzfig{decay_isobar}
\figcaption{fig:decay-isobar}{Measured atomic masses of the isobars with $A=101$, relative to the lightest.
The dashed curve is the best-fitting parabola. Its single minimum is at ruthenium-101, which is indeed the
only stable nuclide with $A=101$. Neutron-rich isobars decay by $\beta^-$ towards it, proton-rich ones by
$\beta^+$ or electron capture. Each step releases the mass difference as energy, so the steps far from the
minimum release the most. Rhodium-101 lies only $\SI{0.54}{\MeV}$ above ruthenium-101, below the
$\SI{1.022}{\MeV}$ threshold, and decays by electron capture alone.}
\end{figure}

Table~\ref{tab:decay-moves} summarises how each mode moves a nucleus on the chart of
Figure~\ref{fig:decay-nzchart}.

\begin{tablebox}{How each decay moves a nucleus on the chart}
\centering
\begin{tabular}{@{}lcccl@{}}
\toprule
\textbf{Decay} & $\Delta Z$ & $\Delta(A-Z)$ & $\Delta A$ & \textbf{Move on the chart}\\
\midrule
$\alpha$ & $-2$ & $-2$ & $-4$ & two steps diagonally down and to the left\\
$\beta^-$ & $+1$ & $-1$ & $0$ & one step towards the valley from above\\
$\beta^+$ or EC & $-1$ & $+1$ & $0$ & one step towards the valley from below\\
$\gamma$ & $0$ & $0$ & $0$ & no move (energy only)\\
\bottomrule
\end{tabular}\label{tab:decay-moves}
\end{tablebox}

\section{Decay Series}\label{sec:decay-series}

A heavy nucleus rarely reaches the valley in one step. An $\alpha$ decay moves it diagonally, parallel to
the line $A-Z=Z$, while the valley bends away from that line; each $\alpha$ therefore leaves the daughter a
little too neutron-rich, and $\beta^-$ decays follow. Uranium-238 passes through fourteen decays before
ending as stable lead-206 (Figure~\ref{fig:decay-useries}). Each $\alpha$ decay lowers $A$ by $4$ and $\beta$
decays leave $A$ unchanged, so every member of a chain has the same remainder when $A$ is divided by $4$.
There are therefore exactly four natural chains:
\begin{itemize}
\item $4n$, the thorium series: $\iso{232}{Th}\to\iso{208}{Pb}$;
\item $4n+1$, the neptunium series: $\iso{237}{Np}\to\iso{209}{Bi}$ (bismuth-209 itself $\alpha$-decays, with a
      half-life of about $2\times10^{19}$ years, to thallium-205). It is extinct in nature: its
      longest-lived member, neptunium-237, has a half-life of $2.1$ million years, and the Earth's
      $4.5$ billion years is over two thousand of its half-lives;
\item $4n+2$, the uranium series: $\iso{238}{U}\to\iso{206}{Pb}$;
\item $4n+3$, the actinium series: $\iso{235}{U}\to\iso{207}{Pb}$.
\end{itemize}

\begin{figure}[htbp]
\centering
\tikzfig{decay_useries}
\figcaption{fig:decay-useries}{The uranium-238 ($4n+2$) series, plotted by mass number $A$ against proton
number $Z$. Diagonal arrows are $\alpha$ decays ($A$ down by $4$, $Z$ down by $2$); curved arrows are
$\beta^-$ decays ($A$ unchanged, $Z$ up by $1$). Every member has $A=4n+2$. Half-lives range from
$4.47$ billion years to $164$ microseconds. The chain contains $8$ $\alpha$ and $6$ $\beta^-$ decays.
Radon-222, a gas, seeps from uranium-bearing rock and is the largest natural source of radiation dose to
the public.}
\end{figure}

\begin{derivbox}{Counting $\alpha$ and $\beta$ decays in a chain}
Go from $\nuc{A}{Z}{X}$ to $\nuc{A'}{Z'}{Y}$ through $n_\alpha$ alpha decays and $n_\beta$ beta-minus
decays, in any order. Only $\alpha$ decay changes $A$, by $4$ each time, so
\[
n_\alpha=\frac{A-A'}{4}.
\]
Each $\alpha$ lowers $Z$ by $2$ and each $\beta^-$ raises it by $1$: $Z'=Z-2n_\alpha+n_\beta$, so
\[
n_\beta=Z'-Z+2n_\alpha.
\]
For $\iso{238}{U}\to\iso{206}{Pb}$: $n_\alpha=32/4=8$ and $n_\beta=82-92+16=6$, as
Figure~\ref{fig:decay-useries} shows. The counts do not depend on the order, which is why they can be found
without knowing the chain.
\end{derivbox}

\begin{exbox}{The first three steps of the uranium series}\label{ex:decay-ufirst}
Uranium-238 undergoes $\alpha$ decay, and the daughter undergoes two successive $\beta^-$ decays. Identify the
three products.

Balance the upper and lower numbers at each step:
$\nuc{238}{92}{U}\xrightarrow{\alpha}\nuc{234}{90}{Th}\xrightarrow{\beta^-}\nuc{234}{91}{Pa}\xrightarrow{\beta^-}\nuc{234}{92}{U}$.
The final product is an \emph{isotope} of the original: one $\alpha$ plus two $\beta^-$ restores $Z$ and
lowers $A$ by $4$. (Uranium-234 makes up $0.005\%$ of natural uranium for exactly this reason.)
\ans{thorium-234, protactinium-234, uranium-234}
\end{exbox}

\begin{exbox}{The thorium series}\label{ex:decay-thseries}
How many $\alpha$ and $\beta^-$ decays take $\nuc{232}{90}{Th}$ to $\nuc{208}{82}{Pb}$?

$n_\alpha=(232-208)/4=6$; $n_\beta=82-90+2(6)=4$. \ans{$6$ $\alpha$, $4$ $\beta^-$}
\end{exbox}

The arithmetic of chains, with its equal steps of $4$ in $A$, says nothing about \emph{how long} each step
takes; Figure~\ref{fig:decay-useries} shows half-lives spanning twenty-one orders of magnitude. How a
population of nuclei decays in time, and how a chain of such populations settles into equilibrium, is the
subject of \chref{ch:decaylaw}.

\section{Summary}\label{sec:summary-decay}
An unstable nucleus moves towards the valley of stability by emitting an $\alpha$ particle, by converting
a neutron into a proton or the reverse, or by shedding excess energy as a photon. Whether a decay is
possible is decided by its $Q$-value; how the energy is shared is decided by momentum conservation.

\begin{itemize}
\item Particles: proton, neutron, electron, positron ($e^+$), neutrino ($\nu_e$), antineutrino
      ($\bar\nu_e$), photon. An antiparticle has the same mass and opposite charges. Lepton number is
      conserved: $e^-$ comes with $\bar\nu_e$, $e^+$ (or a captured $e^-$) with $\nu_e$.
\item Electron--positron annihilation at rest gives two back-to-back $\SI{511}{\keV}$ photons; pair
      production needs a photon above $\SI{1.022}{\MeV}$ and a nearby nucleus.
\item Every decay conserves charge, nucleon number, energy, momentum and lepton number, and needs
      $Q=(m_{\text{initial}}-m_{\text{final}})c^2>0$.
\item $\alpha$: $Z\to Z-2$, $A\to A-4$; $Q_\alpha=[M(\mathrm{X})-M(\mathrm{Y})-M(\iso4{He})]c^2$;
      $K_\alpha=Q(A-4)/A$, $K_{\mathrm{Y}}=4Q/A$; a sharp line.
\item $\beta^-$: $Z\to Z+1$; $Q=[M(\mathrm{X})-M(\mathrm{Y})]c^2$. $\beta^+$: $Z\to Z-1$;
      $Q=[M(\mathrm{X})-M(\mathrm{Y})-2m_e]c^2$. EC: $Z\to Z-1$; $Q=[M(\mathrm{X})-M(\mathrm{Y})]c^2$.
      Atomic masses $M$ throughout.
\item Beta spectra are continuous from $0$ to the endpoint $Q$ (three bodies); the mean electron energy
      is roughly $Q/3$, and the neutrino's share escapes every absorber.
\item $\gamma$: $Z$ and $A$ unchanged; sharp lines from transitions between nuclear levels.
\item Penetration: $\alpha$ (paper) $<$ $\beta$ (mm of aluminium) $<$ $\gamma$ (attenuated by lead, never
      stopped). In a magnetic field, $r=p/|q|B$.
\item At fixed $A$ the atomic mass is a parabola in $Z$; $\beta$ decays run down it to the valley.
\item In a decay series $A$ mod $4$ is fixed; $n_\alpha=\Delta A/4$, $n_\beta=Z'-Z+2n_\alpha$.
\end{itemize}

\section{Problem Types}\label{sec:types-decay}
Seven types cover the decay questions examiners set. Most of them turn on two habits: write every decay
with its upper and lower numbers balanced, and decide before calculating whether the electron masses
cancel.

\archq{Checking a process against the conservation laws}
{``Is this decay or reaction possible?'', ``which particle is missing?'', annihilation or pair
production.}
{Tally, before and after: (1) electric charge; (2) nucleon number; (3) lepton number $L_e$ ($e^-,\nu_e$: $+1$;
$e^+,\bar\nu_e$: $-1$; everything else $0$); (4) energy: $Q>0$ for a spontaneous process; (5) momentum:
in particular, a process with zero total momentum cannot end as a single photon, and a single photon
cannot turn into a pair in empty space. Any failure forbids the process.}
{Checking only charge and nucleon number. Many forbidden processes balance both and fail on lepton number
or momentum.}
{Which of the following processes is forbidden?}
{$\nuc{22}{11}{Na}\to\nuc{22}{10}{Ne}+e^++\nu_e$}
{$n\to p+e^-+\bar\nu_e$}
{$\gamma\to e^-+e^+$ in empty space}
{$\nuc74{Be}+e^-\to\nuc73{Li}+\nu_e$\archans{C}{A lone photon carries momentum $E/c$ in every frame, but in
the frame where the created pair has zero total momentum the photon would need zero momentum; energy and
momentum cannot both be conserved without a nucleus nearby. Options (A), (B) and (D) balance charge,
nucleon number and lepton number, have $Q>0$, and are all observed ($\beta^+$ decay, free-neutron decay
and electron capture).}}

\archq{Writing and completing decay equations}
{``Identify X'', ``complete the equation'', or a chain of named emissions.}
{Balance upper numbers ($A$) and lower numbers (charge) separately: $\alpha=\nuc42{He}$,
$\beta^-=\nuc{0}{-1}{e}$, $\beta^+=\nuc{0}{+1}{e}$; a captured electron counts $-1$ on the left. Read the
element off $Z$. Add the correct neutrino: antineutrino with $e^-$, neutrino with $e^+$ or EC.}
{Pairing the electron with a neutrino instead of an antineutrino, which violates lepton number.}
{Phosphorus-32 ($Z=15$) decays to sulphur-32 ($Z=16$). The particles emitted are:}
{$e^+$ and $\nu_e$}
{$e^-$ and $\nu_e$}
{$e^+$ and $\bar\nu_e$}
{$e^-$ and $\bar\nu_e$\archans{D}{$Z$ rises by $1$ at fixed $A$, so a neutron became a proton: $\beta^-$
decay, which emits an electron ($L_e=+1$) and therefore an antineutrino ($L_e=-1$). Option (B) balances charge
but gives total $L_e=+2$; (A) and (C) would lower $Z$, not raise it, and (C) also fails on lepton
number.}}

\archq{Counting decays in a radioactive series}
{Initial and final nuclides of a chain are given; find the number of $\alpha$ and $\beta$ decays, or decide
whether two nuclides can lie in the same series.}
{$n_\alpha=\Delta A/4$, then $n_\beta=Z_{\text{final}}-Z_{\text{initial}}+2n_\alpha$. Two nuclides share a
series only if their mass numbers leave the same remainder on division by $4$.}
{Forgetting that each $\alpha$ lowers $Z$ by $2$, not $1$, or recalling the familiar counts for a whole
series when the question stops part-way.}
{The numbers of $\alpha$ and $\beta^-$ decays in the chain $\nuc{238}{92}{U}\to\nuc{214}{82}{Pb}$ are:}
{$6$ $\alpha$ and $2$ $\beta^-$}
{$6$ $\alpha$ and $4$ $\beta^-$}
{$8$ $\alpha$ and $6$ $\beta^-$}
{$5$ $\alpha$ and $2$ $\beta^-$\archans{A}{$n_\alpha=(238-214)/4=6$ and $n_\beta=82-92+2(6)=2$. Option (C) is
the count for the whole series down to lead-206; (B) lets each $\alpha$ lower $Z$ by only $1$,
giving $|82-92+6|=4$; (D) would end at $A=218$, not $214$.}}

\archq{$Q$-value of $\alpha$ decay and energy sharing}
{Atomic masses of parent, daughter and $\iso4{He}$ are given; asked for $Q$, the $\alpha$ kinetic energy, the
recoil energy, or a speed ratio.}
{$Q=[M(\mathrm{X})-M(\mathrm{Y})-M(\iso4{He})]\times\SI{931.5}{\MeV}/\si{\amu}$. Then $K_\alpha=Q(A-4)/A$ and
$K_{\mathrm{Y}}=4Q/A$ (equal momenta). Speeds are inversely proportional to masses. Working backwards:
$Q=K_\alpha A/(A-4)$.}
{Giving the products equal energies, or equal speeds, instead of equal momenta; or assuming the heavy
daughter does not move at all.}
{Polonium-210 at rest emits an $\alpha$ particle; the $Q$-value is $\SI{5.41}{\MeV}$. The kinetic energy of
the recoiling lead-206 nucleus is closest to:}
{$\SI{5.30}{\MeV}$}
{$\SI{0.103}{\MeV}$}
{zero}
{$\SI{2.70}{\MeV}$\archans{B}{$K_{\mathrm{Pb}}=4Q/A=4(5.41)/210=\SI{0.103}{\MeV}$. Option (A) is the $\alpha$
particle's share, $Q(206/210)$, the result of giving the products equal speeds; (C) ignores momentum
conservation; (D) shares $Q$ equally.}}

\archq{$\beta^\pm$ and electron-capture $Q$-values}
{A $\beta$ decay with atomic masses; ``is this decay possible?''; endpoint energies.}
{$\beta^-$ and EC: $Q=[M(\mathrm{X})-M(\mathrm{Y})]c^2$. $\beta^+$: subtract $2m_ec^2=\SI{1.022}{\MeV}$. A
decay is possible only if $Q>0$. If gamma rays follow, the $\beta$ endpoint is $Q$ minus the excitation
energy of the level reached.}
{Treating $\beta^+$ like $\beta^-$, or subtracting only one electron mass.}
{The atomic masses of carbon-11 and boron-11 are $\SI{11.011434}{\amu}$ and $\SI{11.009305}{\amu}$. The
$Q$-value for the $\beta^+$ decay of carbon-11 is:}
{$\SI{1.983}{\MeV}$}
{$\SI{1.472}{\MeV}$}
{$\SI{0.961}{\MeV}$}
{$\SI{3.005}{\MeV}$\archans{C}{The atomic mass difference is $\SI{0.002129}{\amu}=\SI{1.983}{\MeV}$, and
for $\beta^+$ two electron masses come off: $1.983-1.022=\SI{0.961}{\MeV}$. Option (A) forgets the
$2m_ec^2$ (it is $Q_{\mathrm{EC}}$); (B) subtracts only one $m_ec^2$; (D) adds $2m_ec^2$ instead of
subtracting it.}}

\archq{Energy sharing in three-body $\beta$ decay}
{An electron of stated energy is observed; asked for the antineutrino's energy, the endpoint, or why the
spectrum is continuous.}
{Neglecting the heavy daughter's recoil, $K_e+E_\nu=Q$, and the endpoint is $K_{\max}=Q$. Two-body decays
give fixed energies; three-body decays give a spectrum.}
{Assuming every electron has the endpoint energy, or that the electron always takes a fixed fraction of
$Q$.}
{Phosphorus-32 decays by $\beta^-$ emission with endpoint energy $\SI{1.71}{\MeV}$. In one decay the electron
has kinetic energy $\SI{0.50}{\MeV}$. Neglecting the recoil of the daughter, the energy carried by the
antineutrino is:}
{$\SI{1.21}{\MeV}$}
{$\SI{1.71}{\MeV}$}
{$\SI{0.50}{\MeV}$}
{$\SI{0.57}{\MeV}$\archans{A}{The endpoint equals $Q$, so $E_{\bar\nu}=Q-K_e=1.71-0.50=\SI{1.21}{\MeV}$.
Option (B) gives the antineutrino everything; (C) assumes an equal split; (D) is $Q/3$, the typical
\emph{average} electron energy, which says nothing about this particular decay.}}

\archq{Decay products in magnetic fields}
{Radii, curvature or direction of tracks in a magnetic field.}
{The magnetic force $|q|vB$ supplies the centripetal force $mv^2/r$, so $r=mv/|q|B=p/|q|B$, with
$p=\sqrt{2mK}$ for a slow particle. For the two products of a decay at rest the momenta are equal, so
$r\propto1/|q|$. The direction follows $q\vv\times\B$.}
{Using $r\propto\sqrt m/q$ (true for equal kinetic energies) when the momenta are equal, or using the mass
ratio.}
{A uranium-238 nucleus at rest in a uniform magnetic field undergoes $\alpha$ decay to thorium-234 ($Z=90$).
Taking both products to carry their full nuclear charges, the ratio $r_\alpha/r_{\mathrm{Th}}$ of the radii
of their circular paths is:}
{$1$}
{$1/45$}
{$58.5$}
{$45$\archans{D}{Equal momenta, so $r\propto1/|q|$: $r_\alpha/r_{\mathrm{Th}}=90/2=45$. Option (A) forgets the
charges; (B) inverts the ratio; (C) is the mass ratio $234/4$, which would apply only to equal speeds.}}

\section{Practice Problems}\label{sec:practice-decay}

The problems below follow the order of the problem types. Atomic masses not given in a problem are those
quoted in the chapter; take $\SI{1}{\amu}\times c^2=\SI{931.494}{\MeV}$ and $m_ec^2=\SI{0.511}{\MeV}$.
Answers, with brief working, follow the list.

\begingroup\emergencystretch=1.5em
\begin{enumerate}[itemsep=3pt]
\item\label{prob:decay-allowed} Which of these are forbidden, and by which law? (a) $n\to p+e^-$;
      (b) $n\to p+e^-+\bar\nu_e$; (c) $p\to n+e^++\bar\nu_e$; (d) $\nuc{14}{6}{C}\to\nuc{14}{7}{N}+e^++\nu_e$;
      (e) an electron and a positron at rest annihilate into one photon.
\item\label{prob:decay-annih} An electron and a positron annihilate at rest. Find the energy of each photon
      and its wavelength ($hc=\SI{1240}{\eV\nano\metre}$).
\item\label{prob:decay-pair} What is the minimum photon energy for pair production? Explain why a photon
      cannot turn into an electron--positron pair in empty space, with no nucleus nearby.
\item\label{prob:decay-na22total} When sodium-22 decays by $\beta^+$ inside a solid, the positron soon
      annihilates. Show that the total energy finally released, counting the neutrino's share, equals
      $Q_{\beta^+}+2m_ec^2$, and compare it with $Q_{\mathrm{EC}}$ for the same nucleus
      (Example~\ref{ex:decay-na22}). How much of it can heat the solid?
\item\label{prob:decay-compl} Complete each equation and name the decay mode:
      (a) $\nuc{22}{11}{Na}\to{?}+e^++\nu_e$; (b) $\nuc{40}{19}{K}+e^-\to{?}+\nu_e$;
      (c) $\nuc{11}{6}{C}\to\nuc{11}{5}{B}+{?}+{?}$; (d) $\nuc{32}{15}{P}\to\nuc{32}{16}{S}+{?}+{?}$.
\item\label{prob:decay-ra2a} Radium-226 ($Z=88$) undergoes two successive $\alpha$ decays. Identify the final
      nuclide.
\item\label{prob:decay-isobar} Which single decay mode turns a nucleus into (a) an isobar, (b) a different
      energy state of the same nuclide? Can any single decay produce an isotope of the parent?
\item\label{prob:decay-u235} Find the numbers of $\alpha$ and $\beta^-$ decays in
      (a) $\nuc{235}{92}{U}\to\nuc{207}{82}{Pb}$; (b) $\nuc{237}{93}{Np}\to\nuc{209}{83}{Bi}$.
\item\label{prob:decay-mod4} To which natural series does $\iso{226}{Ra}$ belong? Can $\iso{232}{Th}$ ever
      decay into $\iso{206}{Pb}$? Explain.
\item\label{prob:decay-back} A chain contains $5$ $\alpha$ decays and $2$ $\beta^-$ decays and ends at
      $\nuc{208}{82}{Pb}$. Find the starting nuclide's $A$ and $Z$.
\item\label{prob:decay-u238a} Given $M(\iso{238}{U})=\SI{238.050788}{\amu}$, $M(\iso{234}{Th})=\SI{234.043601}{\amu}$
      and $M(\iso4{He})=\SI{4.002603}{\amu}$, find $Q$ and $K_\alpha$ for the $\alpha$ decay of uranium-238.
\item\label{prob:decay-po212} Polonium-212 emits $\alpha$ particles of kinetic energy $\SI{8.78}{\MeV}$.
      Find the $Q$-value of the decay.
\item\label{prob:decay-vratio} In the $\alpha$ decay of $\iso{226}{Ra}$ at rest, find the ratio of the speed of
      the $\alpha$ particle to that of the recoiling radon nucleus.
\item\label{prob:decay-be7} $M(\iso7{Be})=\SI{7.016929}{\amu}$ and $M(\iso7{Li})=\SI{7.016003}{\amu}$. Show that
      beryllium-7 cannot decay by positron emission but can decay by electron capture, and find
      $Q_{\mathrm{EC}}$.
\item\label{prob:decay-c14q} Find the $Q$-value for the $\beta^-$ decay of carbon-14, given
      $M(\iso{14}{C})=\SI{14.003242}{\amu}$ and $M(\iso{14}{N})=\SI{14.003074}{\amu}$.
\item\label{prob:decay-trit} Tritium, with atomic mass $M(\iso3{H})=\SI{3.016049}{\amu}$, decays by $\beta^-$
      emission to helium-3, with $M(\iso3{He})=\SI{3.016029}{\amu}$. Find $Q$ in keV. (This tiny
      endpoint makes tritium the best nucleus for measuring the neutrino's mass.)
\item\label{prob:decay-co60} The atomic masses of cobalt-60 and nickel-60 are $\SI{59.933816}{\amu}$ and
      $\SI{59.930786}{\amu}$. Find $Q$ for the $\beta^-$ decay of cobalt-60. The decay feeds a nickel level at $\SI{2.506}{\MeV}$ (Figure~\ref{fig:decay-co60}). Find the
      maximum kinetic energy of the emitted electrons.
\item\label{prob:decay-nde} A free neutron decays ($Q=\SI{0.782}{\MeV}$) and the electron is observed with
      kinetic energy $\SI{0.300}{\MeV}$. Estimate the antineutrino's energy.
\item\label{prob:decay-twobody} Show from momentum and energy conservation that a nucleus at rest decaying
      into exactly \emph{two} particles must give each a unique kinetic energy. Explain why this proved that
      beta decay produces a third particle.
\item\label{prob:decay-mean} The mean kinetic energy of $\beta^-$ electrons is roughly $Q/3$. For the decay of
      cobalt-60 to the $\SI{2.506}{\MeV}$ level of nickel-60 (endpoint $\SI{0.316}{\MeV}$), estimate the
      average energy carried away by antineutrinos per decay, and explain why this energy is lost to any
      absorber.
\item\label{prob:decay-rpa} A proton, a deuteron (charge $+e$, mass $2m_p$) and an $\alpha$ particle with
      equal kinetic energies enter the same uniform magnetic field perpendicularly. Find the ratio of the
      radii of their paths.
\item\label{prob:decay-rrn} A $\iso{226}{Ra}$ nucleus at rest in a uniform magnetic field undergoes $\alpha$
      decay. Treating both products as carrying their full nuclear charges ($+2e$ and $+86e$), find the ratio
      of the radii of the $\alpha$ and radon paths.
\item\label{prob:decay-dirs} A beam containing $\alpha$, $\beta^-$ and $\gamma$ radiation travels in the $+x$
      direction through a magnetic field along $-z$. In which direction is each component initially
      deflected?
\end{enumerate}
\endgroup

\subsection*{Answers}
{\small
\begin{enumerate}[itemsep=2pt]
\item (a) Forbidden by lepton number, $L_e$: $0\to+1$ (and the electron energy would be fixed, contrary to
      observation). (b) Allowed: free-neutron decay, $Q=\SI{0.782}{\MeV}$. (c) Forbidden by lepton number, $L_e$:
      $0\to-1-1=-2$: it would need $\nu_e$, not $\bar\nu_e$ (and a free proton is too light in any case).
      (d) Forbidden by charge, $6\to7+1=8$. (e) Forbidden by momentum: the pair has zero total momentum but
      a single photon always carries momentum $E/c$; two photons are needed.
\item Each photon carries $m_ec^2=\SI{0.511}{\MeV}$. Wavelength
      $=hc/E=\SI{1240}{\eV\nano\metre}/\SI{511000}{\eV}=\SI{2.43e-3}{nm}=\SI{2.43}{pm}$.
\item $2m_ec^2=\SI{1.022}{\MeV}$ (ignoring the small recoil energy of the nucleus). In empty space, view the
      process in the frame where the created pair has zero total momentum: there the photon would need zero
      momentum, but a photon always carries momentum $E/c$, so energy and momentum cannot both be
      conserved. A nearby nucleus absorbs the excess momentum.
\item $Q_{\beta^+}=\SI{1.821}{\MeV}$ is shared between the positron's kinetic energy and the neutrino;
      annihilation then adds $2m_ec^2=\SI{1.022}{\MeV}$. Total $1.821+1.022=\SI{2.843}{\MeV}$, exactly
      $Q_{\mathrm{EC}}$: either route releases the full atomic mass difference
      $[M(\mathrm{X})-M(\mathrm{Y})]c^2$. Only the neutrino's share escapes; the rest (the positron's kinetic
      energy and, if absorbed, the two $\SI{511}{\keV}$ photons) heats the solid.
\item (a) $\nuc{22}{10}{Ne}$, $\beta^+$. (b) $\nuc{40}{18}{Ar}$, electron capture. (c) $e^+$ and $\nu_e$, $\beta^+$.
      (d) $e^-$ and $\bar\nu_e$, $\beta^-$.
\item $\nuc{226}{88}{Ra}\to\nuc{222}{86}{Rn}\to\nuc{218}{84}{Po}$: polonium-218.
\item (a) $\beta^-$, $\beta^+$ or EC ($A$ unchanged). (b) $\gamma$ decay. No single decay gives an isotope:
      $\alpha$ changes $Z$ by $2$ and $\beta$ by $1$.
\item (a) $n_\alpha=28/4=7$, $n_\beta=82-92+14=4$. (b) $n_\alpha=28/4=7$, $n_\beta=83-93+14=4$.
\item $226=4(56)+2$: the $4n+2$ (uranium) series. $232$ is of the form $4n$ but $206$ is $4n+2$; no
      sequence of decays changes the remainder of $A$ on division by $4$, so no.
\item $A=208+5(4)=228$; $82=Z-2(5)+2$ gives $Z=90$: thorium-228.
\item $Q=[238.050788-234.043601-4.002603](931.494)=\SI{4.27}{\MeV}$; $K_\alpha=4.27(234/238)=\SI{4.20}{\MeV}$.
\item $Q=K_\alpha A/(A-4)=8.78(212/208)=\SI{8.95}{\MeV}$.
\item Equal momenta, so $v_\alpha/v_{\mathrm{Rn}}=m_{\mathrm{Rn}}/m_\alpha\approx222/4=55.5$.
\item $\Delta M=\SI{0.000926}{\amu}$, equivalent to $\SI{0.862}{\MeV}<\SI{1.022}{\MeV}$, so
      $Q_{\beta^+}=0.862-1.022=\SI{-0.160}{\MeV}<0$ (forbidden) while $Q_{\mathrm{EC}}=\SI{0.862}{\MeV}>0$.
\item $(14.003242-14.003074)(931.494)=\SI{0.156}{\MeV}$.
\item $(3.016049-3.016029)(931.494)=\SI{0.0186}{\MeV}\approx\SI{18.6}{\keV}$. (The difference is only
      $\SI{0.000020}{\amu}$, so every digit of the masses matters.)
\item $Q=(0.003030)(931.494)=\SI{2.822}{\MeV}$; $K_{\max}=2.822-2.506=\SI{0.316}{\MeV}$.
\item $E_{\bar\nu}\approx0.782-0.300=\SI{0.48}{\MeV}$.
\item At rest the total momentum is zero, so $\vp_1=-\vp_2$ and both particles have the same momentum
      magnitude $p$. Energy: $p^2/2m_1+p^2/2m_2=Q$, so $p^2=2Qm_1m_2/(m_1+m_2)$, a single value, and each
      kinetic energy $K_i=p^2/2m_i$ is fixed. With only two bodies a continuous electron spectrum is
      impossible, so a third, unseen particle must share the energy.
\item Endpoint $\SI{0.316}{\MeV}$; mean electron energy $\approx\SI{0.105}{\MeV}$, so the antineutrinos carry
      on average $\approx\SI{0.21}{\MeV}$ per decay. Neutrinos interact so weakly that they escape any
      absorber, so this energy never appears as heat.
\item $r=\sqrt{2mK}/|q|B\propto\sqrt m/|q|$: $\sqrt1/1:\sqrt2/1:\sqrt4/2=1:\sqrt2:1$.
\item Equal momenta, so $r\propto1/|q|$: $r_\alpha/r_{\mathrm{Rn}}=86/2=43$.
\item $\vF=q\vv\times\B$ with $\uhat{x}\times(-\uhat{z})=+\uhat{y}$. $\alpha$: towards $+y$; $\beta^-$: towards
      $-y$, and usually much more sharply, since its momentum is far smaller; $\gamma$: undeflected.
\end{enumerate}
}
